\documentclass[11pt]{amsart}

\usepackage[T1]{fontenc}
\usepackage[utf8]{inputenc}
\usepackage{lmodern}
\usepackage{amssymb}
\usepackage{mathtools}
\usepackage{booktabs}
\usepackage[margin=1in]{geometry}
\usepackage[hidelinks]{hyperref}

\setlength{\emergencystretch}{3em}

\newcommand{\lean}[1]{\texttt{\detokenize{#1}}}
\newcommand{\fn}[1]{\mathrm{#1}}
\newcommand{\leanref}[1]{\par\nopagebreak\smallskip
  {\small\raggedright\noindent Lean: #1\par}\medskip}

\newtheorem{theorem}{Theorem}
\theoremstyle{definition}
\newtheorem{definition}{Definition}

\title[Pseudosquares primality: a verified proof]{Deterministic primality for sieve survivors by the pseudosquares test:\\
a formally verified proof and a checked table}
\author{Artur Flamandzki}
\date{}

\begin{document}

\begin{abstract}
A partial sieve removes every number with a prime factor up to a bound $s$; its survivors are
then tested for primality. We replace the usual probabilistic test by the pseudosquares test
of Lukes, Patterson and Williams with Sorenson's sieve bound, and verify the whole argument in
Lean~4 / Mathlib. We prove the bound in the form we use (an odd $n > 1$ with every prime factor
above $s$, $n < L s$, that passes the pseudosquares conditions for the bases up to $p$ is a
prime power, where $L$ bounds the pseudosquares for $p$ from below), its consequence for sieve
survivors (with $n < (s+1)^4$, a survivor that passes and is neither a square nor a cube is
prime), and that the test never rejects a prime. The same two properties are proved for a model
of the C routine used in the PrimeAtlas engine, with the table of pseudosquares $L_p$,
$p \le 373$, as its only assumption. The transcribed table is checked by the Lean kernel:
every entry is $\equiv 1 \pmod 8$ and a quadratic residue of every odd prime up to $p$, except
three entries, which are typographical errors that also appear in a published table; each erroneous value lies
below the true one, so the test remained sound. Every statement in this paper is a theorem of
the accompanying Lean formalization, which builds with no \lean{sorry} and uses only the axioms
\lean{propext}, \lean{Quot.sound} and \lean{Classical.choice}.
\end{abstract}

\maketitle

\section{Introduction}

The PrimeAtlas engine lists primes in ranges far above $2^{64}$. One of its paths, \emph{pt1},
sieves a range with the primes up to a bound $s$ and then tests each survivor. Every prime factor
of a survivor exceeds $s$, and $s$ is chosen with $n < (s+1)^4$ for every $n$ in the range, so a
survivor has at most three prime factors. A probabilistic test such as BPSW leaves open whether
the reported numbers are prime. This paper replaces it by a test whose correctness is proved.

The test is the pseudosquares test. The pseudosquare $L_p$ is the least non-square
$N \equiv 1 \pmod 8$ that is a nonzero quadratic residue of every odd prime $q \le p$. Lukes,
Patterson and Williams~\cite{lpw} showed that a number passing certain Euler-criterion conditions
for the bases up to $p$ and lying below $L_p$ is a prime power; Sorenson~\cite{sorenson2006}
extended the bound to $n < L_p s$ for numbers with no prime factor up to $s$, which is exactly the
situation of sieve survivors. The values $L_p$ are known from large computations; the table used
here is the one of~\cite{sorenson2010} for $p \le 373$.

\subsection*{Contributions}
\begin{enumerate}
  \item[(C1)] A formal proof of the pseudosquares bound with the sieve bound $s$
    (Theorem~\ref{thm:lpw}) and of its consequence for sieve survivors (Theorem~\ref{thm:prime}),
    with the table of pseudosquares as an explicit hypothesis.
  \item[(C2)] A formal proof that the test is complete: it never rejects a prime
    (Theorem~\ref{thm:complete}).
  \item[(C3)] Soundness and completeness of a model of the C routine of the engine, including
    its use of a prefix of the bases and its tail loop (Theorem~\ref{thm:routine}).
  \item[(C4)] A kernel check of the transcribed table (Theorem~\ref{thm:table}), which found three
    typographical errors.
\end{enumerate}

\subsection*{Related work}
Primality has been formalized mainly through certificates for individual primes: Pocklington
certificates in Coq~\cite{coqprime}, Pratt certificates in Isabelle~\cite{pratt}, and
Pocklington certificates in Lean~4 in the PrimeCert library~\cite{primecert}. These prove that a
given number is prime. The present work instead proves a test that decides every sieve survivor
of a range, without a certificate for each number. We found no earlier formalization of the
pseudosquares test or of the bound of~\cite{lpw, sorenson2006}; this is a statement about our
search, not a claim of priority.

\subsection*{Notation}
All numbers are natural numbers. $q^{e} \bmod n$ is the remainder of $q^e$ on division by $n$;
``$\equiv -1$'' is written $= n - 1$, as in the formalization. A number $N$ is a
\emph{nonzero square modulo $q$} if $q \nmid N$ and $N$ is a square in $\mathbb{Z}/q$.

\section{Setting}

\begin{definition}[Lower bound for the pseudosquares]
For $p \ge 2$, a number $L$ \emph{bounds the pseudosquares for $p$ from below} if every
$N < L$ with $N \equiv 1 \pmod 8$ that is a nonzero square modulo every odd prime $q \le p$ is a
square. The pseudosquare $L_p$ is the largest such $L$.
\end{definition}

\begin{definition}[The pseudosquares conditions]
For odd $n$ and $p \ge 2$:
\begin{enumerate}
  \item[(P1)] for every prime $q \le p$, $q^{(n-1)/2} \bmod n \in \{1, n-1\}$;
  \item[(P2)] if $n \equiv 5 \pmod 8$, then $2^{(n-1)/2} \bmod n = n - 1$;
  \item[(P3)] if $n \equiv 1 \pmod 8$, then $q^{(n-1)/2} \bmod n = n - 1$ for some prime
    $q \le p$.
\end{enumerate}
\end{definition}

\section{Results}\label{sec:results}

\subsection*{T1 The pseudosquares bound with a sieve bound}
\begin{theorem}\label{thm:lpw}
Let $n > 1$ be odd, let every prime factor of $n$ exceed $s$, let $n < L s$, where $L$ bounds
the pseudosquares for $p \ge 2$ from below, and let $n$ satisfy (P1)--(P3). Then $n = r^k$ for a
prime $r$.
\end{theorem}
\leanref{\lean{lpw_sorenson}}

The proof follows~\cite[Lemma~3.1]{sorenson2006}. Write $n - 1 = 2^e m$ with $m$ odd. For a prime
factor $x$ of $n$ with $x - 1 = 2^f t$, $t$ odd, conditions (P1) and (P3) give $f \ge e$, and the
Legendre symbol $(q \mid x)$ is $-1$ exactly when $q^{(n-1)/2} \bmod n = n - 1$ and $f = e$
(\lean{factor_char}). If some factor has $f > e$, it is a nonzero square modulo every odd
$q \le p$ and $\equiv 1 \pmod 8$, yet it lies below $L$: impossible (\lean{caseA}). If two
distinct factors $r, d$ have $f = e$, they have the same quadratic character, so $r d$ is a
non-square pseudosquare candidate below $L$ (\lean{caseB}). Hence $n$ has one prime factor.

\subsection*{T2 Sieve survivors}
\begin{theorem}\label{thm:prime}
Under the hypotheses of Theorem~\ref{thm:lpw}, if moreover $n < (s+1)^4$ and $n$ is neither a
square nor a cube, then $n$ is prime.
\end{theorem}
\leanref{\lean{prime_of_pss}}

\subsection*{T3 Completeness}
\begin{theorem}\label{thm:complete}
Let $n$ be a prime, $n \ne 2$, $p < n$, and let $L$ bound the pseudosquares for $p$ from below.
Then $n$ satisfies (P1) and (P2), and if $n \equiv 1 \pmod 8$ and $q^{(n-1)/2} \bmod n = 1$ for
every prime $q \le p$, then $L \le n$.
\end{theorem}
\leanref{\lean{pss_complete}}

\subsection*{T4 The C routine}
The table has $74$ rows $(p_k, L_k)$, $k = 0, \dots, 73$, with $p_0 = 2, \dots, p_{73} = 373$.
The routine $\fn{pssProve}(n, b)$ (result $1$: prime, $0$: composite, $-1$: undecided) checks
(P1) for the first $b$ bases, then (P2), accepts when (P3) holds among the first $b$ bases or
$n \not\equiv 1 \pmod 8$, rejects when $n < L_{b-1}$, and otherwise continues base by base
through the rest of the table (\lean{tailProve}). Montgomery arithmetic is modelled as
$x^e \bmod n$.

\begin{theorem}\label{thm:routine}
Assume the table is correct, that is, $L_k$ bounds the pseudosquares for $p_k$ from below for
every $k$ (\lean{TableCorrect}). Let $1 \le b \le 74$.
\begin{enumerate}
  \item If $n > 1$ is odd, every prime factor of $n$ exceeds $s$, $n < (s+1)^4$,
    $n < L_{b-1} s$, $n$ is neither a square nor a cube, and $\fn{pssProve}(n, b) = 1$, then $n$
    is prime.
  \item If $n$ is a prime with $n > 373$, then $\fn{pssProve}(n, b) \ne 0$.
\end{enumerate}
\end{theorem}
\leanref{\lean{pssProve_sound}, \lean{pssProve_complete}}

\subsection*{T5 The table}
\begin{theorem}\label{thm:table}
The bases $p_0, \dots, p_{73}$ are the $74$ primes up to $373$ in increasing order, and the
$L_k$ are non-decreasing. For every $k$ except $k = 49, 70, 71$ ($p = 229, 353, 359$),
$L_k \equiv 1 \pmod 8$ and $(L_k \bmod p_j)^{(p_j - 1)/2} \bmod p_j = 1$ for every
$1 \le j \le k$. The three exceptions are
\[
  L_{49} = 196640148121928601, \qquad L_{70} = L_{71} = 295363187400900310880401;
\]
the first is divisible by $3$; each is below the corrected value
\[
  196640248121928601, \qquad 295363487400900310880401,
\]
and the corrected values pass the same check.
\end{theorem}
\leanref{\lean{tp_all_primes}, \lean{tp_succ}, \lean{tL_succ}, \lean{table_pseudosquares},
\lean{table_typos}, \lean{table_corrected}}

The erroneous digits are a $1$ in place of a $2$ in $L_{229}$ and a $1$ in place of a $4$ in
$L_{353} = L_{359}$. They appear in Table~1 of~\cite{sorenson2010} (in the arXiv version; we found no erratum);
the corrected values agree with the OEIS table~\cite{oeis}, and the corrected value of
$L_{353}$ also appears in~\cite{sorenson2006}. The same three values occur in the table of the
Pseudosquares Prime Sieve~\cite{walisch}. Because each
erroneous value is smaller than the true one, a routine using the erroneous table accepts a subset
of the numbers it would accept with the true table, so its soundness is not affected.

\section{Use in the engine}

On an AMD Ryzen~9 5900X (24 threads, windows of $10^7$), the deterministic path produced the
same prime counts as the BPSW path and as the independent Pseudosquares Prime Sieve of
Walisch~\cite{walisch} on every measured floor from $10^{25}$ to $10^{32}$, and as the engine's
segmented sieve wherever that was run (up to $10^{30}$). The proof costs about nine times the BPSW test and is the whole cost of
the path: $13$--$21$ million numbers per second. These measurements are not part of the formal
results. With the table up to $L_{373}$ the path covers floors up to $10^{34}$ (at $10^{33}$ and
$10^{34}$ with a larger sieve bound); floor $10^{35}$ needs pseudosquares beyond $L_{373}$.

\section{The formalization}\label{sec:formalization}

The formalization is a Lean~4~\cite{lean4} project built against Mathlib~\cite{mathlib}, both at
version \lean{v4.35.0}. All results are in \lean{PtProof/PseudoSquares.lean}; the table checks
are evaluated by the kernel with \lean{decide}. \lean{lake build} completes with no \lean{sorry},
\lean{admit} or \lean{native_decide}, and the theorems depend only on the axioms
\lean{propext}, \lean{Quot.sound} and \lean{Classical.choice}. To build:
\begin{verbatim}
lake exe cache get
lake build
\end{verbatim}

\section{Use of AI}

The decision to make the pt1 path deterministic, the choice of the pseudosquares test and the
statements to be proved are the author's. Generative AI (Claude Opus 5.5, Anthropic, used through
Claude Code) contributed in the following ways:
\begin{itemize}
  \item It wrote the Lean statements and proofs. The author checked the statements against the
    engine; the correctness of the proofs rests on the Lean kernel.
  \item It drafted this text. The author set its structure and the stated contributions.
  \item It searched the literature and checked the identifiers of the references. The author
    decided how the work is placed relative to that literature.
\end{itemize}

\begin{thebibliography}{99}

\bibitem{coqprime}
B.~Gr\'egoire, L.~Th\'ery and B.~Werner,
\emph{A computational approach to Pocklington certificates in type theory},
Functional and Logic Programming (FLOPS 2006), Lecture Notes in Computer Science 3945, 2006, 97--113,
\href{https://doi.org/10.1007/11737414_8}{doi:10.1007/11737414\_8}.

\bibitem{lpw}
R.~F.~Lukes, C.~D.~Patterson and H.~C.~Williams,
\emph{Some results on pseudosquares},
Math. Comp. 65 (1996), 361--372,
\href{https://doi.org/10.1090/S0025-5718-96-00678-3}{doi:10.1090/S0025-5718-96-00678-3}.

\bibitem{mathlib}
The mathlib Community,
\emph{The Lean mathematical library},
Proceedings of the 9th ACM SIGPLAN International Conference on Certified Programs and Proofs
(CPP 2020), 367--381,
\href{https://doi.org/10.1145/3372885.3373824}{doi:10.1145/3372885.3373824}.

\bibitem{lean4}
L.~de~Moura and S.~Ullrich,
\emph{The Lean 4 theorem prover and programming language},
Automated Deduction -- CADE 28, Lecture Notes in Computer Science, 2021, 625--635,
\href{https://doi.org/10.1007/978-3-030-79876-5_37}{doi:10.1007/978-3-030-79876-5\_37}.

\bibitem{oeis}
OEIS Foundation Inc.,
\emph{The On-Line Encyclopedia of Integer Sequences}, sequence A002189 (pseudosquares),
\url{https://oeis.org/A002189}.

\bibitem{pratt}
S.~Wimmer and L.~Noschinski,
\emph{Pratt's primality certificates},
Archive of Formal Proofs, July 2013, ISSN 2150-914x,
\url{https://isa-afp.org/entries/Pratt_Certificate.html}.

\bibitem{primecert}
B.~Mehta et al.,
\emph{PrimeCert: primality certificates in Lean~4},
\url{https://reservoir.lean-lang.org/@b-mehta/PrimeCert}.

\bibitem{sorenson2006}
J.~P.~Sorenson,
\emph{The pseudosquares prime sieve},
Algorithmic Number Theory (ANTS-VII), Lecture Notes in Computer Science 4076, 2006, 193--207,
\href{https://doi.org/10.1007/11792086_15}{doi:10.1007/11792086\_15}.

\bibitem{sorenson2010}
J.~P.~Sorenson,
\emph{Sieving for pseudosquares and pseudocubes in parallel using doubly-focused enumeration and
wheel datastructures},
Algorithmic Number Theory (ANTS-IX), Lecture Notes in Computer Science 6197, 2010, 331--339,
\href{https://doi.org/10.1007/978-3-642-14518-6_26}{doi:10.1007/978-3-642-14518-6\_26};
\href{https://arxiv.org/abs/1001.3316}{arXiv:1001.3316}.

\bibitem{walisch}
K.~Walisch,
\emph{Pseudosquares Prime Sieve},
\url{https://github.com/kimwalisch/Pseudosquares-Prime-Sieve}.

\end{thebibliography}

\end{document}
